\chapter{Systemmodelle}
\label{ch:systemmodelle}

\section{Zielsetzung des Kühlungsmodells}
\label{sec:modell_ziel}

Zur zukünftigen Auslegung und Validierung der Fahrzeugkühlung wird ein
Simulink-Modell des Gesamtsystems aufgebaut. Das Modell bildet alle für die
Kühlung relevanten Komponenten ab bzw. berücksichtigt zumindest deren
thermische Einflussgrößen. Es soll im weiteren Projektverlauf sowohl zur
Dimensionierung neuer Kühlkonzepte als auch zur Validierung bestehender
Konzepte eingesetzt werden.

% FIGURE-VORSCHLAG 1: Übersichts-Blockschaltbild des Gesamtmodells
% (Datenquellen -> Datenaufbereitung -> Thermonetzwerk -> Ausgaben/Bewertung).
% Passt hier ans Ende der Section, als Orientierung für das ganze Kapitel.
\begin{figure}[htbp]
    \centering
    % \includegraphics[width=\textwidth]{bilder/modell_uebersicht}
    \caption{Übersicht über die Struktur des Kühlungsmodells}
    \label{fig:modell_uebersicht}
\end{figure}

\section{Modellaufbau}
\label{sec:modellaufbau}

Der Aufbau des Modells gliedert sich in die folgenden Teilmodule:

% TODO: Liste der Teilmodule einfügen
\begin{itemize}
    \item Akkumulator
    \item Bremse
    \item Motoren
    \item Inverter
\end{itemize}

% FIGURE-VORSCHLAG 4: Simulink-Screenshot der obersten Modellebene
% mit den Teilmodulen.

\subsection{Akkumodell}
\label{subsec:akkumodell}

Das Modell für die Verlustleistung des Akkus wurde bewusst einfach gehalten.
Bei Pouchzellen treten die typischen Hotspots an den Zelltabs auf, insbesondere
am negativen Zelltab~\cite{TODO_Quelle}. Der Zellkörper erwärmt sich ebenfalls,
jedoch deutlich schwächer als die dünnen Tabs. Entsprechend beschränkte sich die
Kühlung des Akkus bisher auf den negativen Zelltab.

% FIGURE-VORSCHLAG 5: Skizze oder Foto einer Pouchzelle mit markierten Zelltabs
% (positiv/negativ) und Hotspot-Bereich, ggf. mit qualitativer Temperaturverteilung.
\begin{figure}[htbp]
    \centering
    % \includegraphics[width=0.7\textwidth]{bilder/pouchzelle_hotspots}
    \caption{Pouchzelle mit Zelltabs und typischem Hotspot am negativen Tab}
    \label{fig:pouchzelle_hotspots}
\end{figure}

Für die Modellierung wurde der Battery Builder aus Simscape Battery in MATLAB
verwendet. Damit lassen sich einzelne Zellen, Stacks und der gesamte Akkuaufbau
abbilden und neben den elektrischen Akkueffekten auch thermisch modellieren.
Der Ansatz erlaubt es, Logging- und Rundenzeitsimulationsdaten als Eingang zu
nutzen, um die Leistungsdaten des Akkus vorzugeben.

Neben den thermischen Simulationen kann das Modell auch zur Akkuauslegung
eingesetzt werden. Diese ist nicht Gegenstand der vorliegenden Arbeit, bietet
aber eine gute Grundlage für eine einfache zukünftige Auslegung von Akku und
Akkukühlung.

% FIGURE-VORSCHLAG 6: Screenshot des Battery Builders bzw. des Akku-Aufbaus
% (Zelle -> Stack -> Pack) mit der verwendeten Konfiguration.
\begin{figure}[htbp]
    \centering
    % \includegraphics[width=0.8\textwidth]{bilder/battery_builder_aufbau}
    \caption{Aufbau des Akkumodells im Battery Builder}
    \label{fig:battery_builder}
\end{figure}

% FIGURE-VORSCHLAG 7 (optional): Simulierte Tab-Temperatur über der Zeit für ein
% Szenario, z.B. Autocross, mit Eingangsleistung des Akkus darüber.

\subsubsection{Bremsen}
\subsubsection{Motoren}
\label{subsubsec:motoren}

Für die Modellierung der Motoren sollte, wie in Abschnitt~\ref{TODO_Abschnitt}
beschrieben, ein Wirkungsgradkennfeld über Messungen erstellt werden. Dies war
nicht möglich. Stattdessen wird das vom Hersteller mitgelieferte
Wirkungsgradkennfeld verwendet.

Das Modell besteht aus einer einfachen thermischen Masse. Eingangsgrößen sind
Drehmoment $M$ und Drehzahl $n$. Aus dem Kennfeld wird der Wirkungsgrad
$\eta(M,n)$ bestimmt und daraus die in den Motor eingebrachte Verlustleistung
berechnet:

\begin{equation}
    P_\mathrm{V} = P_\mathrm{mech} \left( \frac{1}{\eta(M,n)} - 1 \right)
                 = 2\pi\, n\, M \left( \frac{1}{\eta(M,n)} - 1 \right)
    \label{eq:verlustleistung_motor}
\end{equation}

Die Wärmeabfuhr über die Wasserkühlung der Motoren wird ebenfalls
berücksichtigt.

% FIGURE-VORSCHLAG 8: Wirkungsgradkennfeld des Herstellers (Drehmoment über
% Drehzahl, Wirkungsgrad als Konturplot), ggf. mit eingezeichneten
% Betriebspunkten aus einer Messfahrt.
\begin{figure}[htbp]
    \centering
    % \includegraphics[width=0.8\textwidth]{bilder/motor_wirkungsgradkennfeld}
    \caption{Wirkungsgradkennfeld des Motors (Herstellerangabe)}
    \label{fig:motor_kennfeld}
\end{figure}

Die Modellparameter, z.\,B.\ Wärmekapazität und Wärmeübergänge, lassen sich
nur aus anderen Versuchen ableiten oder müssen an vielen Stellen geschätzt
werden. Eine ausführliche Parameteridentifikation steht daher noch aus.

Das Modell ist als einfaches Simscape-Modell aufgebaut. Jeder Motor bildet ein
eigenes Subsystem. Dadurch lassen sich unterschiedliche Aufbauten und
Kühllayouts sowie die Pumpenauswahl validieren.

% FIGURE-VORSCHLAG 9: Simscape-Screenshot eines Motor-Subsystems
% (thermische Masse, Verlustleistungsquelle, Kühlkreis-Anbindung)
% oder alternativ ein thermisches Ersatzschaltbild als Skizze.
\begin{figure}[htbp]
    \centering
    % \includegraphics[width=0.8\textwidth]{bilder/motor_subsystem}
    \caption{Motor-Subsystem im Simscape-Modell}
    \label{fig:motor_subsystem}
\end{figure}

\subsubsection{Inverter}
Kein plan mehr. Ich hab den nie sauber modelliert. Die sind arsch komplex

